% =============================================================
% formatting.tex — Mise en forme imposée par le guide pédagogique
% (Rapport Stage Pédagogique, ESPRIT, §"Forme du Rapport")
% =============================================================

% --- Police : Times New Roman, 12 pt pour le corps du texte ---
\setmainfont{Times New Roman}
% Si "Times New Roman" n'est pas trouvée par XeLaTeX (machine sans
% cette police système), remplacer temporairement par :
%   \setmainfont{Liberation Serif}   (clone métrique libre de Times)

% --- Police à chasse fixe (extraits de config, \ttfamily) : Consolas
%     couvre les caractères de dessin de boîte (─│├└) utilisés dans les
%     arborescences de fichiers, absents de Latin Modern Mono par défaut ---
\setmonofont{Consolas}

% --- Interligne : 1.15 ---
\setstretch{1.15}
\Urlmuskip=0mu plus 1mu\relax

% --- Paragraphes : alinéa de première ligne 0,50 cm et respiration légère
%     après le paragraphe, conformément au gabarit retenu. ---
\setlength{\parindent}{0.5cm}
% L'alinéa de première ligne marque déjà le début de paragraphe : le blanc
% inter-paragraphe est ramené de 2 pt à 0 pt (réserve d'étirement conservée
% pour la mise en colonne). Sur les 133 pages numérotées, ce seul réglage
% rendait quatre pages de blanc sans rien retirer au texte (mesuré 30/08/2026).
\setlength{\parskip}{0pt plus 1pt}

% --- Aucun mot coupé en fin de ligne ---
% Le rapport privilégie les mots entiers, y compris pour les composés. Une
% tolérance modérée et une réserve d'étirement évitent les débordements sans
% transformer le texte en drapeau.
\hyphenpenalty=10000
\exhyphenpenalty=10000
\pretolerance=1500
\tolerance=3500
\emergencystretch=2em

% --- Justification à droite et à gauche (comportement par défaut
%     de LaTeX en mode "flush" — aucune commande supplémentaire). La langue
%     reste le français, mais les pénalités ci-dessus neutralisent la césure. ---

% --- Contrôle des veuves et orphelines (norme ESPRIT §A.2) ---
\clubpenalty=10000
\widowpenalty=10000
\displaywidowpenalty=10000
\raggedbottom

% --- Titres : gras, taille décroissante selon la profondeur (18/16/14/12 pt),
%     numérotés 1, 1.1, 1.1.1 ; retraits normalisés (norme ESPRIT §B.2) :
%       - 0,20 cm de retrait du titre par rapport au bord de la page
%       - 0,30 cm de retrait des lignes suivantes d'un titre sur deux lignes
%       - 0,20 cm entre le numéro et le texte du titre
\titleformat{\chapter}[hang]
  {\hangindent=0.3cm\hangafter=1\normalfont\bfseries\fontsize{18}{22}\selectfont}
  {\chaptertitlename~\thechapter}{0.2cm}{}
\titlespacing*{\chapter}{0.2cm}{0pt}{14pt}

\titleformat{\section}[hang]
  {\hangindent=0.3cm\hangafter=1\normalfont\bfseries\fontsize{16}{19}\selectfont}
  {\thesection}{0.2cm}{}
\titlespacing*{\section}{0.2cm}{10pt}{4pt}

\titleformat{\subsection}[hang]
  {\hangindent=0.3cm\hangafter=1\normalfont\bfseries\fontsize{14}{17}\selectfont}
  {\thesubsection}{0.2cm}{}
\titlespacing*{\subsection}{0.2cm}{7pt}{3pt}

\titleformat{\subsubsection}[hang]
  {\hangindent=0.3cm\hangafter=1\normalfont\bfseries\fontsize{12}{14}\selectfont}
  {\thesubsubsection}{0.2cm}{}
\titlespacing*{\subsubsection}{0.2cm}{5pt}{2pt}

% --- Paragraphes solidaires : un titre ne doit jamais rester seul en
%     bas de page (norme ESPRIT §A.2). Réserve la place d'au moins
%     quelques lignes de texte après chaque titre de section. Équivalent
%     manuel de \Needspace* (paquet "needspace" indisponible hors ligne),
%     construit avec \pagegoal/\pagetotal, primitives du noyau LaTeX. ---
\newcommand{\reservespace}[1]{%
  \par\nobreak
  % Garde 1 : sur une page encore vierge, \pagegoal vaut \maxdimen ;
  % la soustraction n'a alors aucun sens et il n'y a rien a reserver.
  \ifdim\pagegoal=\maxdimen\else
    \dimen0=\pagegoal \advance\dimen0 by -\pagetotal
    \ifdim\dimen0<#1
      % Garde 2 : ne jamais ejecter une page a peine commencee, sous peine
      % de produire une page de trois lignes suivie d'un grand blanc.
      \ifdim\pagetotal>0.4\textheight \newpage \fi
    \fi
  \fi
}
% \pretocmd{\section}{\reservespace{4\baselineskip}}{}{}   % desactive : voir la note ci-dessous
% \pretocmd{\subsection}{\reservespace{3\baselineskip}}{}{}   % desactive : voir la note ci-dessous
% \pretocmd{\subsubsection}{\reservespace{3\baselineskip}}{}{}   % desactive : voir la note ci-dessous
% NOTE (26/08/2026) : les trois accrochages ci-dessus sont desactives.
% La macro lit la position courante AVANT que TeX ait coupe le paragraphe en
% cours : la valeur lue est celle de la page precedente, la garde passe, et le
% saut de page produit exactement l'artefact qu'elle pretend eviter, soit une
% page de trois lignes suivie d'un grand blanc. Les penalites de veuve et
% d'orpheline, fixees plus haut a 10000, suffisent a empecher un titre isole en
% bas de page. La macro est conservee pour un eventuel appel manuel.

% --- Sous-titre de chapitre (italique, centré, non numéroté, absent de
%     la table des matières) — norme ESPRIT §B.3 / plan détaillé ---
\newcommand{\soustitrechapitre}[1]{%
  \begingroup
  \centering\itshape\large #1\par
  \endgroup
  \vspace{10pt}%
}

% --- Profondeur de numérotation et de table des matières : chapitre,
%     section, sous-section (norme ESPRIT §B.1, recommandation de ne
%     pas descendre en dessous du niveau 3) ---
\setcounter{secnumdepth}{3}
\setcounter{tocdepth}{1}

% --- Tableaux longs : une taille légèrement réduite améliore la lecture des
%     matrices denses et limite les ruptures, sans descendre sous 8 pt. ---
\AtBeginEnvironment{longtable}{\scriptsize}

% --- Grille des tableaux : bordure extérieure et séparation de chaque ligne
%     et de chaque colonne. La largeur réduite des marges internes compense les
%     filets verticaux sans dépasser la largeur utile de la page. ---
\setlength{\tabcolsep}{3pt}
\setlength{\arrayrulewidth}{0.4pt}
\renewcommand{\arraystretch}{1.0}

% --- Placement des flottants ---
% Les valeurs par défaut de LaTeX sont conservatrices : une page portant une
% figure doit conserver au moins 20 % de texte (\textfraction) et un flottant
% de tête ne peut occuper plus de 70 % de la hauteur (\topfraction). Avec une
% trentaine de figures, dont plusieurs occupent la moitié d'une page, la
% conjonction des deux laissait des pages remplies à 55 % suivies d'un grand
% blanc. Les seuils sont donc relâchés : une figure peut désormais partager sa
% page avec le texte qui la suit, ce qui densifie le document sans rien retirer.
% Blancs autour des flottants et de leurs légendes. Les valeurs de la classe
% (20 pt autour d'un flottant, 12 pt sous une légende) ont été calibrées pour
% des documents à cinq ou six illustrations ; le mémoire en compte plus de
% quatre-vingts, et ce blanc seul coûtait environ deux pages.
\setlength{\textfloatsep}{8pt plus 3pt minus 2pt}
\setlength{\intextsep}{8pt plus 3pt minus 2pt}
\setlength{\floatsep}{8pt plus 3pt minus 2pt}
\setlength{\abovecaptionskip}{4pt}
\setlength{\belowcaptionskip}{0pt}
\setlength{\LTpre}{6pt}
\setlength{\LTpost}{6pt}

\renewcommand{\topfraction}{0.92}
\renewcommand{\bottomfraction}{0.70}
\renewcommand{\textfraction}{0.06}
\renewcommand{\floatpagefraction}{0.72}
\setcounter{topnumber}{3}
\setcounter{bottomnumber}{2}
\setcounter{totalnumber}{4}

% --- Numérotation des pages : bas à droite, "page / total" ---
% Confirmé par la norme ESPRIT §A.1 : « en bas à droite, sous la forme
% page/total » — implémenté ci-dessous comme "n / N".
\pagestyle{fancy}
\fancyhf{}
\fancyfoot[R]{\thepage\,/\,\pageref{LastPage}}
\renewcommand{\headrulewidth}{0pt}

% Les pages d'ouverture de chapitre utilisent le style "plain" par
% défaut : on le redéfinit pour garder la pagination bas-droite.
\fancypagestyle{plain}{
  \fancyhf{}
  \fancyfoot[R]{\thepage\,/\,\pageref{LastPage}}
  \renewcommand{\headrulewidth}{0pt}
}

% --- Numérotation continue des figures, tableaux et graphiques (pas de
%     remise à zéro par chapitre), séparée par type (norme ESPRIT §C) :
%     Fig. 1, Tab. 1, Graph. 1, chaque type comptant indépendamment sur
%     l'ensemble du document. Le compteur "graphique" (package float,
%     cf. packages.tex) ne compte pas par chapitre par construction. ---
\counterwithout{figure}{chapter}
\counterwithout{table}{chapter}
\captionsetup[figure]{name=Fig.,labelsep=period}
\captionsetup[table]{name=Tab.,labelsep=period}
\captionsetup[graphique]{labelsep=period}
% --- Francisation de cleveref (defaut L2) ---
% Les \crefname doivent imperativement etre poses dans \AtBeginDocument.
% En preambule, les \renewcommand des conjonctions s'accumulent dans
% \extras<langue> et provoquent un depassement de pile au premier
% \selectlanguage (reproduit et isole le 26/08/2026).
\AtBeginDocument{%
  \crefname{figure}{figure}{figures}%
  \Crefname{figure}{Figure}{Figures}%
  \crefname{table}{tableau}{tableaux}%
  \Crefname{table}{Tableau}{Tableaux}%
  \crefname{section}{section}{sections}%
  \Crefname{section}{Section}{Sections}%
  \crefname{subsection}{section}{sections}%
  \Crefname{subsection}{Section}{Sections}%
  \crefname{subsubsection}{section}{sections}%
  \Crefname{subsubsection}{Section}{Sections}%
  \crefname{chapter}{chapitre}{chapitres}%
  \Crefname{chapter}{Chapitre}{Chapitres}%
  \crefname{graphique}{graphique}{graphiques}%
  \Crefname{graphique}{Graphique}{Graphiques}%
  \renewcommand{\crefpairconjunction}{ et }%
  \renewcommand{\crefmiddleconjunction}{, }%
  \renewcommand{\creflastconjunction}{ et }%
  \renewcommand{\crefrangeconjunction}{ \`a }%
  \renewcommand{\crefpairgroupconjunction}{ et }%
  \renewcommand{\creflastgroupconjunction}{ et }%
}

% --- Recto verso : pas d'en-tête/pied de page sur les pages blanches
%     insérées par \cleardoublepage avant chaque chapitre. Reproduit
%     manuellement l'effet du paquet "emptypage" (indisponible hors
%     ligne), selon la méthode standard documentée. ---
\makeatletter
\let\origcleardoublepage\cleardoublepage
\newcommand{\clearemptydoublepage}{%
  \clearpage
  {\pagestyle{empty}\origcleardoublepage}%
}
\renewcommand{\cleardoublepage}{\clearemptydoublepage}
\makeatother

% --- Liens hypertexte discrets (impression recto-verso conservée) ---
\hypersetup{
  colorlinks=true,
  linkcolor=black,
  citecolor=black,
  urlcolor=blue!50!black,
  pdfborder={0 0 0}
}
